% ==========================================================================
%  Anhang – Quellenmaterial und Arbeitsdokumente aus dem Ordner "ressourcen"
% ==========================================================================

% Lokales Hilfsmakro: Bild/Skizze auf Seitenmass begrenzt einbetten.
\newcommand{\AnhangBild}[2]{%
  \begin{center}%
    \adjustbox{max width=\textwidth,max totalheight=0.80\textheight}{\includegraphics{#1}}\\[0.4em]%
    {\small\itshape #2}%
  \end{center}%
  \par\medskip
}

Dieser Anhang bündelt das Arbeits- und Quellenmaterial, das im Verlauf des
Projekts entstanden ist. Er ist nach den Phasen des Double-Diamond-Prozesses
gegliedert. Abschnitt~\ref{anh:uebersicht} listet sämtliche Dateien des
Ordners \texttt{ressourcen} auf; die folgenden Abschnitte binden die
zentralen Visualisierungen, Skizzen sowie die textbasierten Arbeitsdokumente
direkt ein.

% --------------------------------------------------------------------------
\subsection{Übersicht aller Quellen und Arbeitsdokumente}
\label{anh:uebersicht}

\begin{footnotesize}
\setlength{\LTpre}{0.5em}\setlength{\LTpost}{0.5em}
\begin{longtable}{@{}p{8.2cm} p{14.6cm}@{}}
\textbf{Datei} & \textbf{Inhalt / Beschreibung}\\
\hline
\endfirsthead
\textbf{Datei} & \textbf{Inhalt / Beschreibung}\\
\hline
\endhead
\multicolumn{2}{@{}l}{\textbf{Übergreifend}}\\
\texttt{Bericht Praxis-Case ZVV.docx} & Word-Arbeitsfassung des Gesamtberichts.\\[0.3em]

\multicolumn{2}{@{}l}{\textbf{Einleitung}}\\
\texttt{einleitung/titelbild.png} & Titelbild des Berichts.\\
\texttt{einleitung/symbole.png} & Symbol- und Bildwelt der Persona Kurt.\\[0.3em]

\multicolumn{2}{@{}l}{\textbf{Identify}}\\
\texttt{identify/deconstructionMap.png} & Deconstruction Map des ZVV-Case.\\
\texttt{identify/designSpaceMap.png} & Design Space Map (Wissen / Annahmen / Nichtwissen).\\
\texttt{identify/stakeHolderMap.png} & Stakeholder Map.\\[0.3em]

\multicolumn{2}{@{}l}{\textbf{Discover}}\\
\texttt{discover/empathymap.png} & Empathy Map der frisch Pensionierten.\\
\texttt{discover/paingain.png} & Pain-/Gain-Übersicht.\\
\texttt{discover/Zusammenfassung\_Interviews\_CAS\_AISD.md} & Zusammenfassung der Interview-Theorie (Skript Blume).\\
\texttt{discover/CAS AISD\_Exploration \dots\ 2026.pdf} & Skript \enquote{Exploration der Kundenperspektive} (59 S.).\\
\texttwemoji{1f4c4}~\texttt{Fragebogen zur Freizeitmobilität\_V2.docx} & Fragebogen zur Freizeitmobilität (Version~2).\\
\texttt{01\_interviews/README.md} & Planung und Übersicht der Interviewserie.\\
\texttt{01\_interviews/Interview Bernhard \_Jaya.docx} & Interviewprotokoll Bernhard.\\
\texttt{01\_interviews/Interview Hans Jörg \_ ZHAW-Bereinigt.docx} & Interviewprotokoll Hans Jörg (bereinigt).\\
\texttt{01\_interviews/Interview Margrit \_ ZHAW .docx} & Interviewprotokoll Margrit (Rohfassung).\\
\texttt{01\_interviews/Interview Margrit \_ ZHAW\_Bereinigt.docx} & Interviewprotokoll Margrit (bereinigt).\\
\texttwemoji{1f4c4}~\texttt{Fragebogen \dots\_Rolf Kohler.doc} & Ausgefüllter Fragebogen (Rolf Kohler).\\
\texttwemoji{1f4c4}~\texttt{Fragebogen \dots\_Toni\_Venzin.docx} & Ausgefüllter Fragebogen (Toni Venzin).\\
\texttwemoji{1f4c4}~\texttt{Fragebogen \dots\_V2\_B. Siegenthaler.docx} & Ausgefüllter Fragebogen (B. Siegenthaler).\\
\texttwemoji{1f4c4}~\texttt{Fragebogen \dots\_V2\_P. Büki.docx} & Ausgefüllter Fragebogen (P. Büki).\\
\texttwemoji{1f4c4}~\texttt{Fragebogen \dots\_V2\_P\_Fiktiv.docx} & Fiktiv ergänzter Fragebogen (Persona-Test).\\
\texttwemoji{1f4c4}~\texttt{Fragebogen \dots\_V2\_S. Dias.docx} & Ausgefüllter Fragebogen (S. Dias).\\[0.3em]

\multicolumn{2}{@{}l}{\textbf{Define}}\\
\texttt{define/clustering\_1.png} & Interview-Mapping: kontinuierliche Achsen (Teil~1).\\
\texttt{define/clustering\_2.png} & Interview-Mapping: kontinuierliche Achsen (Teil~2).\\
\texttt{define/clustering\_3.png} & Interview-Mapping: diskrete Zuordnungen (Teil~1).\\
\texttt{define/clustering\_4.png} & Interview-Mapping: diskrete Zuordnungen (Teil~2).\\
\texttt{define/customer\_journey.md} & Customer Journey über 12 Phasen.\\
\texttt{define/HMW\_Kurt\_Gesamtliste.docx} & Gesamtliste der How-Might-We-Fragen.\\
\texttt{define/Ideengenerierung.docx} & Protokoll der Ideengenerierung.\\
\texttt{define/Ideengenerierung\_Claude\_Martina.docx} & Ideengenerierung (KI-gestützt, Martina).\\
\texttt{define/Ideen\_Kurt\_HowNowWow.docx} & Bewertung der Ideen (How-Now-Wow-Matrix).\\
\texttt{define/Ideen\_Masterliste\_Kurt.docx} & Konsolidierte Ideen-Masterliste.\\
\texttt{define/top10\_ideen\_ZVV.docx} & Top-10-Ideen für den ZVV.\\
\texttt{define/Archive/hmw-fragen-kurt.md} & HMW-Fragen nach Clustern (Markdown).\\
\texttt{define/Archive/HmW Fragen.xlsx} & HMW-Sammlung (Excel).\\
\texttt{define/Archive/HMW\_KurtKreuzUndQuer\_ZVV.docx} & HMW-Archiv (Word).\\
\texttt{define/Archive/How Might We Fragen.docx} & HMW-Archiv (Word).\\[0.3em]

\multicolumn{2}{@{}l}{\textbf{Develop}}\\
\texttt{develop/develop1.png} & Brainstorming-Board 1 (Miro).\\
\texttt{develop/develop2.png} & Brainstorming-Board 2 (Miro).\\
\texttt{develop/develop3.png} & Brainstorming-Board 3 (Miro).\\
\texttt{develop/kopfstandmethode.png} & Kopfstand-Methode (Reverse Brainstorming).\\
\texttt{develop/ideenskizzen/Ausflugs\_Baukasten\_Ideenskizze.pdf} & Ideenskizze \enquote{Ausflugs-Baukasten}.\\
\texttt{develop/ideenskizzen/Ausflugs\_Community\_Ideenskizze.pdf} & Ideenskizze \enquote{Ausflugs-Community}.\\
\texttt{develop/ideenskizzen/Ausflugs\_Rezeptbuch\_Ideenskizze.pdf} & Ideenskizze \enquote{Ausflugs-Rezeptbuch}.\\
\texttt{develop/ideenskizzen/Customer\_Journey\_Ideen\_Gaps.pdf} & Ideen-Gaps entlang der Customer Journey.\\
\texttt{develop/ideenskizzen/Fahrt\_Quiz\_Ideenskizze.pdf} & Ideenskizze \enquote{Fahrt-Quiz}.\\
\texttt{develop/ideenskizzen/QR\_3in1\_Ideenskizze.pdf} & Ideenskizze \enquote{QR 3-in-1}.\\
\texttt{develop/ideenskizzen/Proto1\_AusflugsRezeptBuch.jpg} & Prototyp 1: Ausflugs-Rezeptbuch.\\
\texttt{develop/ideenskizzen/Proto2\_AllInOneQrCode.jpg} & Prototyp 2: All-in-One-QR-Code.\\
\texttt{develop/ideenskizzen/Proto3\_PlanerGameification.jpg} & Prototyp 3: Planer mit Gamification.\\
\texttt{develop/ideenskizzen/Proto4\_AusflugsCommunity.jpg} & Prototyp 4: Ausflugs-Community.\\
\texttt{develop/ideenskizzen/Proto5\_TagesreiseApp.jpg} & Prototyp 5: Tagesreise-App.\\
\texttt{develop/ideenskizzen/Proto6\_DonnerstagsTreffpunkt.jpg} & Prototyp 6: Donnerstags-Treffpunkt.\\
\texttt{develop/CAS\_AISD\_Slides\_Ideation\_Konzeption\_FS2026.pdf} & Vorlesungsslides Ideation \& Konzeption (71 S.).\\
\texttt{develop/Gruppe 2 Simona.pptx} & Präsentationsfoliensatz Gruppe~2.\\
\texttt{develop/OneDrive\_\dots/Plakat Donnerstag mit Kurt.png} & Plakatentwurf \enquote{Donnerstag mit Kurt}.\\
\texttt{develop/OneDrive\_\dots/ZVV Plakat.png} & Plakatentwurf ZVV Tagestouren.\\
\texttt{develop/OneDrive\_\dots/ZVV Gruppenreisen\_Midjourney.png} & KI-Bild ZVV Gruppenreisen.\\
\texttt{develop/OneDrive\_\dots/ZVV Gruppenreisen Ende.png} & KI-Bild ZVV Gruppenreisen (Abschluss).\\
\texttt{develop/OneDrive\_\dots/ZVV Kirche.png} & KI-Bild Ausflug (Altstadt / Kirche).\\
\texttt{develop/OneDrive\_\dots/ZVV ZUG.png} & KI-Bild Gruppe im Zug.\\
\texttt{develop/OneDrive\_\dots/ZVV Tagestouren V2.mp4} & Videomaterial Tagestouren (V2).\\
\texttt{develop/OneDrive\_\dots/ZVV Tagestouren V2\_Voice.mp4} & Videomaterial Tagestouren (V2, Voice).\\
\texttt{develop/OneDrive\_\dots/ZVV-Guide-Recruiting-Film.mp4} & Recruiting-Film für ZVV-Guides.\\
\texttt{develop/OneDrive\_\dots/Claude Design Apps.url} & Link zu genutzten Design-Apps.\\[0.3em]

\multicolumn{2}{@{}l}{\textbf{Deliver}}\\
\texttt{deliver/Enrico.png} & Profil der Testperson Enrico (71).\\
\texttt{deliver/Lisbeth.png} & Profil der Testperson Lisbeth (68).\\
\texttt{deliver/idee.md} & Konzeptbeschrieb \enquote{Multiplikator Kurt}.\\
\texttt{deliver/Feedback\_Capture\_Grid\_ZVV\_G2.pdf} & Feedback Capture Grid der Tests (Basis).\\
\texttt{deliver/Feedback\_Capture\_Grid\_ZVV\_G2 (1)\dots(4).pdf} & Feedback Capture Grids der Tests (Varianten 1--4).\\
\texttt{deliver/Interview-Leitfaden.docx} & Leitfaden für die Testinterviews.\\
\texttt{deliver/Testing mit Enrico \_\_ ZHAW G2.docx} & Testprotokoll mit Enrico.\\
\texttt{deliver/Testing mit Lisbeth \_\_ ZHAW G2.docx} & Testprotokoll mit Lisbeth.\\
\end{longtable}
\end{footnotesize}

% ==========================================================================
\clearpage
\pagecolor{anhangPastel}
\subsection{Identify}

\AnhangBild{ressourcen/identify/deconstructionMap.png}{Deconstruction Map: Zerlegung des ZVV-Case in Zielgruppe, Rahmenbedingungen und Stossrichtungen.}
\clearpage
\AnhangBild{ressourcen/identify/designSpaceMap.png}{Design Space Map: \enquote{Was wir wissen / annehmen / nicht wissen}.}
\clearpage
\AnhangBild{ressourcen/identify/stakeHolderMap.png}{Stakeholder Map mit Kernnutzern, direktem Umfeld und weiteren Akteuren.}

% ==========================================================================
\clearpage
\pagecolor{anhangPastel}
\subsection{Discover}

\AnhangBild{ressourcen/discover/empathymap.png}{Empathy Map der frisch Pensionierten (65--70 J.).}
\clearpage
\AnhangBild{ressourcen/discover/paingain.png}{Pain- und Gain-Übersicht aus den Interviews.}

\clearpage
\subsubsection{Übersicht der Interviewserie}
\textit{Quelle: \texttt{discover/01\_interviews/README.md}}

\medskip
Es sollen mindestens 10 Interviews durchgeführt werden:
\begin{itemize}[nosep]
  \item 2 organisierte Teams-Calls von uns
  \item 6 Calls von anderen Teams
  \item 4 zusätzliche Befragungen durch uns
  \item X Probanden via Martinas Fragebogen
\end{itemize}

\subsubsection{Zusammenfassung: Exploration der Kundenperspektive durch Interviews}
\textit{Quelle: \texttt{discover/Zusammenfassung\_Interviews\_CAS\_AISD.md} (Marc Blume, CAS AISD FS~2026)}

\paragraph{Was ist ein Interview?}
Ein zielgerichtetes Gespräch zur Ermittlung von Informationen: Der Interviewer
stellt Fragen, die befragte Person antwortet in eigenen Worten, die Aussagen
werden meist aufgezeichnet. Kernunterschied zum Alltagsgespräch: strukturiert,
zielgerichtet und rollenbasiert.

\paragraph{Wann eignet sich ein Interview?}
\begin{itemize}[nosep]
  \item \textbf{Zugänglichkeit:} wenn eine Beobachtung nicht möglich ist.
  \item \textbf{Zeithorizont:} wenn die Ereignisse in Vergangenheit oder Zukunft liegen.
  \item \textbf{Perspektive:} wenn die Innensicht der Befragten zentral ist.
  \item \textbf{Voraussetzung:} die Person muss fähig und bereit sein, Auskunft zu geben.
\end{itemize}

\paragraph{Interviewformen}
\begin{enumerate}[nosep]
  \item \textbf{Standardisiert:} Wortlaut und Reihenfolge fest vorgegeben -- für hohe Vergleichbarkeit.
  \item \textbf{Leitfadeninterview (halbstandardisiert):} flexibel, max.\ 2~Seiten mit ca.\ 8--15 Fragen; setzt Vorwissen voraus. \texttwemoji{26a0}~Achtung: oft zu früh eingesetzt -- erst explorieren, dann strukturieren.
  \item \textbf{Narratives Interview (nichtstandardisiert):} freies Erzählen; hoher Freiraum; Phasen Erzählanstoss, Haupterzählung, Nachfragephase.
\end{enumerate}

\paragraph{Vorbereitung in 5 Schritten}
\begin{enumerate}[nosep]
  \item Ziele der Interviewserie bestimmen.
  \item Vorhandenes Vorwissen sichten (Literatur, KI-Recherche, Experten).
  \item Leitfaden entwerfen: Themen listen, Fragen formulieren.
  \item In Pilotphase erproben und zyklisch verbessern.
  \item Interviewserie durchführen und Leitfaden laufend anpassen.
\end{enumerate}

\paragraph{6 Regeln für gute Interviewfragen}
\begin{enumerate}[nosep]
  \item Sprache des Interviewpartners verwenden.
  \item \enquote{Warum?}-Fragen sparsam einsetzen.
  \item Keine Ja/Nein-Fragen.
  \item Immer nur eine Frage auf einmal.
  \item Keine doppelten Verneinungen.
  \item Keine Suggestivfragen.
\end{enumerate}

\paragraph{Aktives Zuhören -- drei Ebenen}
\begin{itemize}[nosep]
  \item \texttwemoji{1f91d}~\textbf{Beziehungsebene:} volle Aufmerksamkeit, offene Körperhaltung, Blickkontakt, keine Unterbrechungen, kurze Bestätigungen.
  \item \texttwemoji{1f3a4}~\textbf{Sachebene:} Aussagen paraphrasieren, Verständnisfragen stellen, Sprechpausen aushalten, aktiv zuhören.
  \item \texttwemoji{2764}~\textbf{Gefühlsebene:} ausgedrückte Gefühle verbalisieren, direkt nachfragen, nonverbale Signale beachten, eigene Meinung zurückstellen.
\end{itemize}

\paragraph{\texttwemoji{26a0}~Zu vermeidende Fehler}
Zu strikt am Leitfaden festhalten, ungeduldig sein, Gesprächspausen nicht
aushalten, Distanz/Intimität schlecht balancieren, die befragte Person
unterbrechen.

\paragraph{Biases vermeiden}
\begin{itemize}[nosep]
  \item \texttwemoji{1f607}~\textbf{Confirmation Bias:} bestätigende Informationen werden bevorzugt -- \enquote{Was nicht passt, wird passend gemacht}.
  \item \texttwemoji{1f440}~\textbf{Halo-Effekt:} ein positives Merkmal strahlt auf die Bewertung anderer Merkmale ab.
  \item \texttwemoji{1f465}~\textbf{Ingroup Bias:} Mitglieder der eigenen Gruppe werden wohlwollender bewertet.
\end{itemize}

\paragraph{Dokumentation}
Die Dokumentationsform muss zum Interviewcharakter passen: offenes
Interview~$\to$~flexible Form (Post-its, Notizen); standardisiertes
Interview~$\to$~strukturierte Form (Word, Excel). Weitere Formen: Papier /
Handnotizen (diskret, aber lückenhaft) und Audioaufnahmen (vollständig, aber
Zustimmung nötig).

\paragraph{Die 5 Key Messages}
\begin{enumerate}[nosep]
  \item Interviews ermöglichen tiefen Einblick -- die Person muss aber Auskunft geben wollen und können.
  \item Je weniger Vorwissen, desto mehr Freiraum für den Interviewpartner.
  \item Ein guter Leitfaden besteht aus entscheidungs- und designrelevanten Fragen -- nicht nur \enquote{interessanten}.
  \item Ein gelungenes Interview liefert Erkenntnisse und fühlt sich angenehm und authentisch an.
  \item Interviewverlauf und Dokumentationsform müssen sich entsprechen.
\end{enumerate}

\paragraph{Zu Beginn des Interviews klären}
Kennenlernen (Name, Rolle, Institution), Zweck und Auswahlgrund erläutern,
Themenüberblick geben, Rahmenbedingungen (Anonymisierung, Datenschutz,
Einwilligung zur Aufzeichnung, Dauer) klären und Rückfragen ermöglichen.

% ==========================================================================
\clearpage
\pagecolor{anhangPastel}
\subsection{Define}

\subsubsection{Interview-Mapping und Clustering}
\label{anh:clustering}
\textit{Quelle: Stimmt-Workshop (Anja Gähwiler), eigene Verortung der Interviewpersonen.}

\medskip
Die folgenden Maps zeigen, wie die Interviewpersonen entlang kontinuierlicher
(Bereich \enquote{sehr tief}--\enquote{sehr hoch}) und diskreter Variablen verortet
und zu Mustern geclustert wurden. Der \gls{oev}-affine Cluster um Guido, Peter
und Hans Jörg (Basis der Leit-Persona Kurt) liegt konsistent am oberen Ende der
ÖV-bezogenen Achsen.

\AnhangBild{ressourcen/define/clustering_1.png}{Interview-Mapping, kontinuierliche Achsen (Teil~1): digitale Affinität, Mobilitätsbedarf im Alltag, ÖV- und Auto-Nutzung, Wohlfühlen und allgemeines Sicherheitsgefühl im ÖV.}
\clearpage
\AnhangBild{ressourcen/define/clustering_2.png}{Interview-Mapping, kontinuierliche Achsen (Teil~2): ÖV-Affinität des Umfelds, Umweltbewusstsein, Lernmotivation für neue Apps, Bedürfnis nach sozialer Interaktion, Grösse/Intaktheit des sozialen Umfelds und Gesundheitsbeschwerden.}
\clearpage
\AnhangBild{ressourcen/define/clustering_3.png}{Interview-Mapping, diskrete Zuordnungen (Teil~1): Motivation für die ÖV-Nutzung, emotionale Funktion des eigenen Autos, Nutzungsgrund für das Auto und Wohnort.}
\clearpage
\AnhangBild{ressourcen/define/clustering_4.png}{Interview-Mapping, diskrete Zuordnungen (Teil~2): Informationsgewinnung, Rolle bei der Freizeitgestaltung, HundehalterIn sowie Grösse und Intaktheit des sozialen Umfelds.}

\clearpage
\subsubsection{Customer Journey (12 Phasen)}
\textit{Quelle: \texttt{define/customer\_journey.md}}

\newcommand{\CJPhase}[6]{%
  \par\smallskip\noindent\textbf{#1}
  \begin{itemize}[nosep,leftmargin=1.4em]
    \item \textbf{Beschreibung:} #2
    \item \textbf{Touchpoints:} #3
    \item \textbf{Pleasure Points:} #4
    \item \textbf{Pain Points:} #5
    \item \textbf{Emotion:} #6
  \end{itemize}%
}

\CJPhase{1 -- Inspiration}
{Idee entsteht spontan aus Eigeninitiative, Tagesanzeiger, Veranstaltungshinweisen oder Vereinsaktivitäten.}
{Tages-Anzeiger, SBB-Newsletter, Vereinsmitteilungen, Wanderwege-App, Mundpropaganda.}
{Breites ÖV-Angebot eröffnet viele Ziele; Vereinsaktivität als starker Anker.}
{Freizeitimpulse kommen selten direkt vom ZVV; wenig zielgruppengerechte Inspiration.}
{\texttwemoji{1f60a}~Neugier, Vorfreude -- ÖV als selbstverständlicher Ermöglicher.}

\CJPhase{2 -- Entscheid}
{Schneller, meist spontaner Entscheid -- kein grosses Abwägen. ÖV ist selbstverständliche Wahl.}
{Persönlicher Kalender, Wetterapps, Intuition.}
{Spontanität möglich -- kein Vorausbuchen nötig; Halbtax/GA senkt die Hemmschwelle.}
{Auto nötig, wenn ÖV nicht hinführt (z.\,B.\ Pfannenstiel, abgelegene Ziele).}
{\texttwemoji{1f604}~Gelassen, souverän -- kein Stress, Entscheidung fällt leicht.}

\CJPhase{3 -- Koordination}
{Kurze Absprache wenn nötig; oft alleine oder mit bekannten ÖV-affinen Personen unterwegs.}
{Telefon/WhatsApp mit Begleitung, eigene Entscheidung bei Alleinreisen.}
{Reisen mit Gleichgesinnten macht Spass; alleine fahren ebenso problemlos möglich.}
{Umfeld fährt oft Auto $\to$ soziale Reibung; Koordination mit Autofahrenden aufwändig.}
{\texttwemoji{1f60a}~Entspannt wenn alleine; leicht frustriert über das Auto-Umfeld.}

\CJPhase{4 -- Ausflug planen}
{Minimalplanung: Zielort, Wanderweg, Wetterprognose und Webcam prüfen.}
{Wanderwege-App, SBB-App, Wetter-App, Webcam, Internet-Suche.}
{Schnelle, zuverlässige Infos via Apps; geringer Planungsaufwand dank Netzkenntnissen.}
{GA gilt nicht überall (Bergbahnen, lokale Netze); Zusatzbillett-Pflicht nicht immer ersichtlich.}
{\texttwemoji{1f60a}~Fokussiert, neugierig -- Planung macht Spass.}

\CJPhase{5 -- Verbindung wählen}
{Aktive Suche nach optimaler Verbindung via SBB-/ZVV-App inkl.\ Umsteigezeiten-Optimierung.}
{SBB-App, ZVV-App, Fahrzielanzeige, Zonenplan (Internet).}
{Freude am Verbindungsoptimieren (\enquote{Gehirnjogging}); schnellere Strecken durch Netzkenntnis.}
{Zonensystem und Anschlussbillett-Regeln schwer verständlich; Upgrade-Regeln verwirrend.}
{\texttwemoji{1f604}~Freude am Optimieren; leichte Irritation bei unklaren Regeln.}

\CJPhase{6 -- Ticketkauf}
{Digital via SBB-App, Halbtax oder GA; Automat als Backup. Kauf läuft routiniert.}
{SBB-App (mobil), Halbtax/GA, Ticketautomat (Karte).}
{Reibungsloser Kauf; Halbtax spart Geld; digitaler Fahrschein ist bequem.}
{Anschlussbillett-Logik komplex; Upgrade muss vor Einsteigen gelöst sein -- viele wissen das nicht.}
{\texttwemoji{1f60a}~Routiniert, entspannt -- nur bei Anschlussbillett-Fragen kurz unsicher.}

\CJPhase{7 -- Weg zum Bahnhof}
{Zu Fuss oder mit dem Velo zum Bahnhof. Kurze, vertraute Wege, kein Stress.}
{Gehweg / Velo, Bahnhofsareal, Anzeigetafeln.}
{Kurzer Fussweg / Velofahrt; Bewegung angenehm; kein Parkplatzstress.}
{Bei Velo: kein ÖV am Ziel (z.\,B.\ Pfannenstielgipfel), daher doch Auto nötig.}
{\texttwemoji{1f604}~Positiv, aktiv -- Bewegung als Genuss.}

\CJPhase{8 -- Fahrt}
{Entspannte Fahrt: Landschaft geniessen, lesen, gelegentlich mit Mitreisenden plaudern.}
{Zug, Sitzplatz, Fensterausblick, Mitreisende, SBB-App für Live-Status.}
{Landschaft, Lesen, Entspannung, Kontakt mit Mitreisenden; keine Fahrkonzentration nötig.}
{Volle Züge (auch ausserhalb HVZ); Stehen müssen; Werbung auf Trams wirkt aggressiv.}
{\texttwemoji{1f604}~Zufrieden, erholt; leicht genervt bei Überfüllung.}

\CJPhase{9 -- Anschluss}
{Kenntnisreiches Umsteigen; kurze Anschlüsse werden bewusst einkalkuliert.}
{Bahnhofsanzeige, SBB-App, Streckenkenntnisse, Bahnhofspersonal.}
{Kurze Anschlüsse kein Problem dank Erfahrung; Pünktlichkeit des Schweizer ÖV.}
{Auslandsanschlüsse unsicher (Italien: Bus kommt nicht); Schiff/Seilbahn-Saisonstart unklar.}
{\texttwemoji{1f60a}~Kompetent, sicher; Stress nur bei unerwarteten Ausfällen.}

\CJPhase{10 -- Erlebnis}
{Wandern, Vereinsaktivität, Kultur oder Stadtbummel -- ÖV ist Mittel zum Zweck.}
{Natur, Museen, Vereinsareal, Restaurants, Bahnangebote vor Ort.}
{Erlebnis selbst ist die Hauptmotivation; ÖV ermöglicht direkte Anreise ohne Parksuche.}
{Teils fehlende ÖV-Anbindung am Zielort zwingt zu Auto oder langen Fusswegen.}
{\texttwemoji{1f604}~Begeistert, erfüllt -- Höhepunkt der Journey.}

\CJPhase{11 -- Rückreise}
{Rückfahrt wird ähnlich geplant. Gedränge am Sonntagabend wird akzeptiert, aber nicht geliebt.}
{SBB-App, Zug, evtl.\ Verspätungsanzeige.}
{Routinemässige Rückreise; bei guter Pünktlichkeit entspannter Abschluss.}
{Gedränge am Sonntagabend nervt; volle Züge sind ein Ärgernis.}
{\texttwemoji{1f60a}~Entspannt bis leicht genervt (Gedränge); müde, aber zufrieden.}

\CJPhase{12 -- Evaluation}
{Mentale Notiz über die Verbindungsqualität; wenig formale Nachbearbeitung.}
{Persönliche Erinnerung, evtl.\ Gespräch mit Begleitung.}
{Zufriedenheit über gute Verbindungen bestätigt die ÖV-Wahl; positives Gesamterlebnis.}
{Keine etablierte Feedback-Möglichkeit; Verbesserungspotenzial \enquote{geht verloren}.}
{\texttwemoji{1f60a}~Dankbar, positiv bestätigt -- ÖV-Affinität wird bekräftigt.}

\clearpage
\subsubsection{How-Might-We-Fragen}
\textit{Quelle: \texttt{define/Archive/hmw-fragen-kurt.md}\quad Basis: Interviews Reini (66), Hans Jörg (78), Peter (66) und Fragebogen-Insights.}

\paragraph{Cluster 1: ÖV als selbstverständliches Default}
\begin{itemize}[nosep]
  \item Wie könnten wir Kurts Überzeugung, dass ÖV einfach Sinn macht, so erlebbar machen, dass sie auf sein Umfeld abstrahlt?
  \item Wie könnten wir Kurts ÖV-Alltag so kuratieren, wie Spotify eine Discover Weekly zusammenstellt -- vertraut, aber mit überraschenden neuen Verbindungen?
  \item Wie könnten wir Kurt eine ÖV-Identität geben, die so sichtbar ist, dass andere fragen: \enquote{Wie bist du dahin gekommen?}
\end{itemize}

\paragraph{Cluster 2: Tracking, Messen \& geistig fit bleiben}
\begin{itemize}[nosep]
  \item Wie könnten wir Kurts Freude am Messen und Dokumentieren mit seinen ÖV-Reisen verbinden, sodass jede Fahrt zum messbaren Erlebnis wird?
  \item Wie könnten wir Kurt in 10 Jahren eine Mobilitätsbiografie aus tausend ÖV-Reisen präsentieren -- und heute damit anfangen?
  \item Wie könnten wir die Tracking-Motivation auch für einen 82-jährigen Kurt ohne Smartphone aktivieren?
\end{itemize}

\paragraph{Cluster 3: Vorbild \& Wissen weitergeben}
\begin{itemize}[nosep]
  \item Wie könnten wir Kurts ÖV-Wissen als Ressource für andere nutzbar machen, ohne dass er dafür extra Aufwand betreiben muss?
  \item Wie könnten wir Kurts ÖV-Kompetenz aus ZVV-Perspektive denken und daraus eine offizielle Rolle gestalten?
  \item Wie könnten wir verhindern, dass Kurts Wissen mit seiner Pensionierung verschwindet -- und es stattdessen zum Startpunkt seiner grössten Wirkung machen?
\end{itemize}

\paragraph{Cluster 4: Frustration über Komplexität \& Ausschluss}
\begin{itemize}[nosep]
  \item Wie könnten wir Kurt den ÖV zugänglich machen, ohne anzunehmen, dass Tarifkomplexität unvermeidlich ist?
  \item Wie könnten wir Kurts ÖV-Erlebnis so gestalten, dass Zonengrenzen und Ticketregeln komplett unsichtbar werden?
  \item Wie könnten wir Kurt ermöglichen, seine Frustration über unlogische Regeln direkt in Systemverbesserungen zu übersetzen?
\end{itemize}

\paragraph{Cluster 5: Geselligkeit \& die Fahrt als Erlebnis}
\begin{itemize}[nosep]
  \item Wie könnten wir eine ÖV-Welt schaffen, in der der Zug nicht Transportmittel, sondern Ort des Erlebnisses ist?
  \item Wie könnten wir Kurt Zugehörigkeit vermitteln, wenn er mit Vereinskollegen eine Ausflugsstrecke plant und einsteigt?
  \item Wie könnten wir die soziale ÖV-Erfahrung auch für ein neu zugezogenes Vereinsmitglied gestalten -- und Kurts Einladungsrolle stärken?
\end{itemize}

% ==========================================================================
\clearpage
\pagecolor{anhangPastel}
\subsection{Develop}

\AnhangBild{ressourcen/develop/kopfstandmethode.png}{Kopfstand-Methode: \enquote{Wie halten wir 65+ von der ÖV fern?} als Ausgangspunkt.}
\clearpage
\AnhangBild{ressourcen/develop/develop1.png}{Brainstorming-Board 1.}
\clearpage
\AnhangBild{ressourcen/develop/develop2.png}{Brainstorming-Board 2.}
\clearpage
\AnhangBild{ressourcen/develop/develop3.png}{Brainstorming-Board 3.}

\clearpage
\subsubsection{Ideenskizzen}
\AnhangBild{ressourcen/develop/ideenskizzen/Ausflugs_Baukasten_Ideenskizze.pdf}{Ideenskizze \enquote{Ausflugs-Baukasten}.}
\clearpage
\AnhangBild{ressourcen/develop/ideenskizzen/Ausflugs_Community_Ideenskizze.pdf}{Ideenskizze \enquote{Ausflugs-Community}.}
\clearpage
\AnhangBild{ressourcen/develop/ideenskizzen/Ausflugs_Rezeptbuch_Ideenskizze.pdf}{Ideenskizze \enquote{Ausflugs-Rezeptbuch}.}
\clearpage
\AnhangBild{ressourcen/develop/ideenskizzen/Customer_Journey_Ideen_Gaps.pdf}{Ideen-Gaps entlang der Customer Journey.}
\clearpage
\AnhangBild{ressourcen/develop/ideenskizzen/Fahrt_Quiz_Ideenskizze.pdf}{Ideenskizze \enquote{Fahrt-Quiz}.}
\clearpage
\AnhangBild{ressourcen/develop/ideenskizzen/QR_3in1_Ideenskizze.pdf}{Ideenskizze \enquote{QR 3-in-1}.}

\clearpage
\subsubsection{Prototypen der Ideen}
\AnhangBild{ressourcen/develop/ideenskizzen/Proto1_AusflugsRezeptBuch.jpg}{Prototyp 1: Ausflugs-Rezeptbuch.}
\clearpage
\AnhangBild{ressourcen/develop/ideenskizzen/Proto2_AllInOneQrCode.jpg}{Prototyp 2: All-in-One-QR-Code.}
\clearpage
\AnhangBild{ressourcen/develop/ideenskizzen/Proto3_PlanerGameification.jpg}{Prototyp 3: Planer mit Gamification.}
\clearpage
\AnhangBild{ressourcen/develop/ideenskizzen/Proto4_AusflugsCommunity.jpg}{Prototyp 4: Ausflugs-Community.}
\clearpage
\AnhangBild{ressourcen/develop/ideenskizzen/Proto5_TagesreiseApp.jpg}{Prototyp 5: Tagesreise-App.}
\clearpage
\AnhangBild{ressourcen/develop/ideenskizzen/Proto6_DonnerstagsTreffpunkt.jpg}{Prototyp 6: Donnerstags-Treffpunkt.}

\clearpage
\subsubsection{Plakat- und Bildmaterial (KI-generiert)}
\AnhangBild{ressourcen/develop/OneDrive_2_16.6.2026/Plakat Donnerstag mit Kurt.png}{Plakatentwurf \enquote{Donnerstag mit Kurt}.}
\clearpage
\AnhangBild{ressourcen/develop/OneDrive_2_16.6.2026/ZVV Plakat.png}{Plakatentwurf ZVV Tagestouren.}
\clearpage
\AnhangBild{ressourcen/develop/OneDrive_2_16.6.2026/ZVV Gruppenreisen_Midjourney.png}{KI-Bild: ZVV Gruppenreisen.}
\clearpage
\AnhangBild{ressourcen/develop/OneDrive_2_16.6.2026/ZVV Gruppenreisen Ende.png}{KI-Bild: ZVV Gruppenreisen (Abschluss).}
\clearpage
\AnhangBild{ressourcen/develop/OneDrive_2_16.6.2026/ZVV Kirche.png}{KI-Bild: Ausflug in der Altstadt.}
\clearpage
\AnhangBild{ressourcen/develop/OneDrive_2_16.6.2026/ZVV ZUG.png}{KI-Bild: Gruppe im Zug.}

% ==========================================================================
\clearpage
\pagecolor{anhangPastel}
\subsection{Deliver}

\subsubsection{Testpersonen}
\AnhangBild{ressourcen/deliver/Enrico.png}{Profil der Testperson Enrico (71).}
\clearpage
\AnhangBild{ressourcen/deliver/Lisbeth.png}{Profil der Testperson Lisbeth (68).}

\clearpage
\subsubsection{Konzeptbeschrieb \enquote{Multiplikator Kurt}}
\textit{Quelle: \texttt{deliver/idee.md}}

\medskip
\textbf{Hauptidee:} Multiplikator Kurt -- ihm das Erlebnis, das Planen, das
Buchen und das Einladen so einfach wie möglich machen, damit der agile,
digital-affine Kurt neue Menschen zum ÖV-Fahren motivieren kann.

\begin{enumerate}[nosep]
  \item \textbf{Die Fahrt ist das Ziel:} Tageserlebnisse/Pakete passend für 65+ schaffen.
  \item \textbf{Unterstützung der Planer:} Planen, Übermitteln, Buchen und Koordinieren der Ausflugsinfos so einfach wie möglich machen.
  \item \textbf{Sicherheits- und Zugehörigkeitsgefühl der Teilnehmenden:} Sorgenfreiheit schaffen -- \enquote{Ich bin dabei und bin umsorgt}.
  \item \textbf{Gamification:} dem Planer (Sammler) den verdienten Status ermöglichen.
  \item \textbf{Belohnung der Kurts:} je nach Menge der Buchungen/Ausflüge fährt Kurt gratis usw.
  \item \textbf{Promotion / Launch-Plan} des Angebots in seiner ganzen Fülle (wie erreichen wir die Kurts dieser Welt).
  \item \textbf{Community-Plattform \enquote{Reis mit mir}:} Anschluss an die Donnerstags-Gruppe. Kampagne mit QR-Code-Registrierung für Ausflüge per Mail. Treffpunkt Zürich HB; die Webseite zeigt die jeweiligen Donnerstags-Tagesausflüge mit Planer (Kurt), Datum, Kurzprofil und Kontaktnummer. Kurt hat das Buchungsportal und löst nach Eintreffen aller Teilnehmenden das Ticket für 8 Personen.
  \item \textbf{Promotionsplan:} Tagesreisen und Gamification zur Akquise von Kurts. Bei genügend Kurts zwei parallele Kampagnen -- eine für Planer, eine für Teilnehmende (\enquote{Schliess dich Kurt an}).
\end{enumerate}

\clearpage
\subsubsection{Feedback Capture Grids der Tests}
Die folgenden Seiten binden die Feedback Capture Grids der Testsessions
vollständig ein (Quelldateien \texttt{deliver/Feedback\_Capture\_Grid\_ZVV\_G2*.pdf}).

\includepdf[pages=-,frame,scale=0.80,pagecommand={\thispagestyle{fancy}}]{ressourcen/deliver/Feedback_Capture_Grid_ZVV_G2.pdf}
\includepdf[pages=-,frame,scale=0.80,pagecommand={\thispagestyle{fancy}}]{ressourcen/deliver/Feedback_Capture_Grid_ZVV_G2 (1).pdf}
\includepdf[pages=-,frame,scale=0.80,pagecommand={\thispagestyle{fancy}}]{ressourcen/deliver/Feedback_Capture_Grid_ZVV_G2 (2).pdf}
\includepdf[pages=-,frame,scale=0.80,pagecommand={\thispagestyle{fancy}}]{ressourcen/deliver/Feedback_Capture_Grid_ZVV_G2 (3).pdf}
\includepdf[pages=-,frame,scale=0.80,pagecommand={\thispagestyle{fancy}}]{ressourcen/deliver/Feedback_Capture_Grid_ZVV_G2 (4).pdf}
